% =============================================================================
% Derivation of the q(beta_k) Variational Update
% Self-Contained Reference for Supervised Bayesian MF Implementation
%
% Author:  Andrew Walther
% Date:    March 2026
%
% PURPOSE:
%   This document derives the coordinate-ascent variational inference (CAVI)
%   update for the survival coefficients beta in the Supervised Bayesian
%   Matrix Factorization model.  The update is cast as an Empirical Bayes
%   Normal Means (EBNM) problem and solved via the ebnm R package.
%
%   This is the first of four modular derivation documents:
%     qBeta   (this document) --> qL --> qF --> qTau
%   Each document covers a single CAVI step and maps directly to a
%   standalone R function (code/update_beta.R, etc.).
%
% COMPANION CODE: code/update_beta.R
% PARENT DOCUMENT: derivations/MF_UpdateDerivations/MF_Derivations_UpdateAlgo_REVISED.tex
% =============================================================================

\documentclass[11pt]{article}
\usepackage[margin=1in]{geometry}
\usepackage{amsmath, amssymb, amsthm}
\usepackage{bm}
\usepackage{xcolor}
\usepackage{hyperref}
\usepackage{booktabs}
\usepackage{array}
\usepackage{cancel}
\usepackage{tcolorbox}
\tcbuselibrary{skins, breakable}

% ---- Same custom commands as MF_Derivations_UpdateAlgo_REVISED.tex ----------
% Posterior mean
\newcommand{\lbar}[2]{\bar{l}_{#1#2}}
\newcommand{\fbar}[2]{\bar{f}_{#1#2}}
\newcommand{\betabar}[1]{\bar{\beta}_{#1}}
% Second moment (E[x^2])
\newcommand{\lsec}[2]{\overline{l^2_{#1#2}}}
\newcommand{\fsec}[2]{\overline{f^2_{#1#2}}}
\newcommand{\betasec}[1]{\overline{\beta^2_{#1}}}
% Partial residuals
\newcommand{\Rk}[0]{R^{-k}}
\newcommand{\zk}[0]{z^{-k}}
% Revision callouts
\newcommand{\revised}[1]{{\color{red}\textbf{[REVISED]:} #1}}

% ---- Same coloured boxes as REVISED.tex -------------------------------------
\newtcolorbox{correctionbox}{
  colback=red!5!white, colframe=red!50!black,
  title=Correction Note, fonttitle=\bfseries\small, breakable
}
\newtcolorbox{ebnmbox}{
  colback=blue!5!white, colframe=blue!40!black,
  title=EBNM Problem, fonttitle=\bfseries\small, breakable
}
\newtcolorbox{derivbox}{
  colback=gray!6!white, colframe=gray!50!black,
  title=Derivation Steps, fonttitle=\bfseries\small, breakable
}
% ---- New box for self-consistency checks ------------------------------------
\newtcolorbox{verifybox}{
  colback=green!5!white, colframe=green!50!black,
  title=Verification Check, fonttitle=\bfseries\small, breakable
}
\newtcolorbox{codebox}{
  colback=yellow!5!white, colframe=orange!50!black,
  title=Code Mapping (\texttt{update\_beta.R}), fonttitle=\bfseries\small, breakable
}

% ---- Document ---------------------------------------------------------------
\title{\textbf{Derivation of the $q_{\beta_k}$ Variational Update}\\[4pt]
       \large Supervised Bayesian Matrix Factorization\\[4pt]
       \normalsize Self-Contained Reference for \texttt{code/update\_beta.R}}
\author{Andrew Walther}
\date{March 2026}

\begin{document}
\maketitle

\begin{abstract}
This document derives the coordinate-ascent variational inference (CAVI)
update for the survival coefficient $\beta_k$ in the Supervised Bayesian
Matrix Factorization (Supervised MF) model.  Because $\beta$ appears only
in the Cox proportional-hazards survival likelihood (not in the genomics
term), the update depends solely on that likelihood.  A second-order Taylor
approximation makes the update tractable, and the resulting quadratic
objective reduces to a one-dimensional Empirical Bayes Normal Means (EBNM)
problem.  The EBNM is solved using the \texttt{ebnm} R package with a
\texttt{point\_normal} prior, which automatically shrinks irrelevant factors
toward zero.

The derivation follows the notation of the parent document
\texttt{MF\_Derivations\_UpdateAlgo\_REVISED.tex} and maps directly to the
standalone R implementation in \texttt{code/update\_beta.R}.  This is the
first of four modular derivation documents (q$\beta$ $\to$ qL $\to$ qF
$\to$ q$\tau$).
\end{abstract}

\tableofcontents
\newpage

% =============================================================================
\section{Setup and Notation}\label{sec:setup}
% =============================================================================

\subsection{Model}

The Supervised MF model combines a genomics matrix factorization with a Cox
proportional-hazards survival model:

\begin{align}
  \text{Genomics:}\quad &
  Y_{ij} = \sum_{k=1}^K l_{ik} f_{jk} + \varepsilon_{ij},
  \quad \varepsilon_{ij} \sim \mathcal{N}(0,\,\tau_j^{-1})
  \label{eq:genomics} \\[4pt]
  \text{Survival:}\quad &
  h(t_i) = h_0(t_i)\exp\!\left(\eta_i\right),
  \quad \eta_i = \sum_{k=1}^K l_{ik}\beta_k = \bm{l}_i^\top\bm{\beta}
  \label{eq:survival}
\end{align}

\noindent
where $Y$ is $n\times p$, $L$ is $n\times K$ (patient loadings),
$F$ is $p\times K$ (biological factors), $\bm{\beta}\in\mathbb{R}^K$
are survival coefficients, $\tau_j$ is the feature-specific noise precision,
and $(t_i, \delta_i)$ are the observed survival time and event indicator for
patient $i$.

\subsection{Key Notation}

Throughout this document, $q(\cdot)$ denotes the variational posterior and
$g(\cdot)$ denotes the (learned) empirical Bayes prior.  The critical
distinction between posterior means and second moments:

\begin{center}
\begin{tabular}{lll}
\toprule
\textbf{Symbol} & \textbf{Definition} & \textbf{R variable} \\
\midrule
$\lbar{i}{k} = \bar{l}_{ik}$ & $\mathbb{E}_q[l_{ik}]$ (posterior mean) & \texttt{EL[i,k]} \\
$\lsec{i}{k} = \overline{l^2_{ik}}$ & $\mathbb{E}_q[l_{ik}^2]$ (posterior 2nd moment) & \texttt{EL2[i,k]} \\
$\betabar{k} = \bar{\beta}_k$ & $\mathbb{E}_q[\beta_k]$ (posterior mean) & \texttt{EBeta[k]} \\
$\betasec{k} = \overline{\beta^2_k}$ & $\mathbb{E}_q[\beta_k^2]$ (posterior 2nd moment) & \texttt{EBeta2[k]} \\
$\zk_i = z_i^{-k}$ & Partial working response (Section~\ref{sec:zk}) & \texttt{z\_no\_k} \\
$W_{ii}$ & Cox neg-diagonal Hessian (positive) & \texttt{w[i]} \\
$z_i$ & Cox working response $\hat{\eta}_i + u_i/W_{ii}$ & \texttt{z[i]} \\
\bottomrule
\end{tabular}
\end{center}

\begin{verifybox}
\textbf{Second moment identity.}  For any random variable $X$ with
finite variance:
\[
  \mathbb{E}_q[X^2] = \mathrm{Var}_q(X) + \bigl(\mathbb{E}_q[X]\bigr)^2
  \;\geq\; \bigl(\mathbb{E}_q[X]\bigr)^2
\]
with equality if and only if the posterior is degenerate (zero variance).
Whenever we write $\overline{l^2_{ik}}$, this refers to the \emph{full}
second moment, not the squared mean.  In code: \texttt{EL2[i,k] = sd\^{}2 +
mean\^{}2} from the EBNM result.
\end{verifybox}

\subsection{What is Fixed During the $\beta$ Update}

When updating $q(\beta_k)$, the following quantities are held fixed:
\begin{itemize}
  \item Posterior moments of loadings: $\bar{l}_{ik}$, $\overline{l^2_{ik}}$ (from the L update of this iteration)
  \item Cox Taylor working quantities: $z_i$, $W_{ii}$ (computed at the start of the outer loop)
  \item All $\bar{\beta}_{k'}$ and $\overline{\beta^2_{k'}}$ for $k' \neq k$ (fixed within the k-loop)
  \item Factor moments $\bar{f}_{jk}$, $\overline{f^2_{jk}}$, and precision $\tau_j$ (from their own updates)
\end{itemize}

% =============================================================================
\section{The Objective Function for $q_{\beta_k}$}\label{sec:objective}
% =============================================================================

The full Evidence Lower Bound (ELBO) to maximise is:
\begin{equation}
  \mathcal{F} = \mathbb{E}_q[\log P(Y, t, \delta \mid L, F, \bm{\beta}, \tau)]
  + \sum_k \mathbb{E}_{q_{l_k}}\!\left[\log\tfrac{g_{l_k}}{q_{l_k}}\right]
  + \sum_k \mathbb{E}_{q_{f_k}}\!\left[\log\tfrac{g_{f_k}}{q_{f_k}}\right]
  + \sum_k \mathbb{E}_{q_{\beta_k}}\!\left[\log\tfrac{g_{\beta_k}}{q_{\beta_k}}\right]
  \label{eq:elbo-full}
\end{equation}

The full data likelihood splits into genomics and survival terms:
\begin{equation}
  \log P(Y, t, \delta \mid L, F, \bm{\beta}, \tau)
  = \underbrace{\log P(Y \mid L, F, \tau)}_{\text{genomics (no }\beta)}
  + \underbrace{\log P(t, \delta \mid L, \bm{\beta})}_{\text{survival}}
  \label{eq:likelihood-split}
\end{equation}

\textbf{Key observation:} $\beta$ does not appear in the genomics term.
Therefore, the genomics log-likelihood is \emph{constant} with respect to
$\beta_k$ and drops out of the objective.  The objective for
$q_{\beta_k}$ is:

\begin{equation}
  \boxed{
  \mathcal{F}(q_{\beta_k}, g_{\beta_k}) =
  \mathbb{E}_{q_l, q_\beta}\!\left[\log P(t, \delta \mid L, \bm{\beta})\right]
  + \sum_k \mathbb{E}_{q_{\beta_k}}\!\left[\log\frac{g_{\beta_k}(\beta_k)}{q_{\beta_k}(\beta_k)}\right]
  }
  \label{eq:objective-beta}
\end{equation}

% =============================================================================
\section{Taylor Approximation of the Cox Likelihood}\label{sec:taylor}
% =============================================================================

The Cox partial log-likelihood $\ell(\bm{\eta})$ is non-conjugate in $L$ and
$\bm{\beta}$.  We apply a second-order Taylor expansion around the current
expansion point $\hat{\eta}_i = \sum_k \bar{l}_{ik}\bar{\beta}_k$:

\begin{equation}
  \ell(\bm{\eta}) \approx \ell(\hat{\bm{\eta}})
  + (\bm{\eta} - \hat{\bm{\eta}})^\top \nabla\ell(\hat{\bm{\eta}})
  + \tfrac{1}{2}(\bm{\eta} - \hat{\bm{\eta}})^\top \nabla^2\ell(\hat{\bm{\eta}}) (\bm{\eta} - \hat{\bm{\eta}})
\end{equation}

where the \textbf{score} $u_i$ and \textbf{neg-diagonal Hessian} $W_{ii}$ are:
\begin{align}
  u_i &= \delta_i - \exp(\hat{\eta}_i) \sum_{j:\,t_j \le t_i}
          \frac{\delta_j}{\sum_{m:\,t_m \ge t_j}\exp(\hat{\eta}_m)}
  \quad\text{(score)} \label{eq:score} \\
  W_{ii} &= \exp(\hat{\eta}_i) \sum_{j:\,t_j \le t_i}
             \frac{\delta_j}{\sum_{m:\,t_m \ge t_j}\exp(\hat{\eta}_m)}
  \quad\text{(neg-diag Hessian, }W_{ii}>0\text{)} \label{eq:hessian}
\end{align}

Completing the square in $\eta_i$ yields the quadratic surrogate:
\begin{equation}
  \log P(t, \delta \mid L, \bm{\beta}) \approx
  \sum_i \left[-\frac{1}{2} W_{ii}\left(\eta_i - z_i\right)^2\right] + C
  \label{eq:taylor-result}
\end{equation}
where the \textbf{working response} is:
\begin{equation}
  z_i = \hat{\eta}_i + \frac{u_i}{W_{ii}}
  \label{eq:working-response}
\end{equation}

The full derivation of this approximation is in the parent document
(\texttt{REVISED.tex}, Section 3).  In code: \texttt{calc\_cox\_taylor()},
lines 70--93 of \texttt{Supervised\_Bayesian\_MF\_V2.R}.

% =============================================================================
\section{Partial Working Response $\zk_i$}\label{sec:zk}
% =============================================================================

To update $\beta_k$ via coordinate ascent, we isolate the dependence on
$\beta_k$ by defining the \textbf{partial working response}:
\begin{equation}
  \zk_i = z_i^{-k} := z_i - \sum_{k' \neq k} \bar{l}_{ik'}\bar{\beta}_{k'}
  \label{eq:zk-def}
\end{equation}

This removes all contributions from other factors, leaving the residual
signal that factor $k$ should explain.

\subsection{Reuse Across L and $\beta$ Updates}

\begin{verifybox}
\textbf{z\_no\_k reuse rationale.}
Note that $z_i^{-k}$ depends only on $k' \neq k$:
\[
  z_i^{-k} = z_i - \underbrace{\sum_{k' \neq k} \bar{l}_{ik'}\bar{\beta}_{k'}}_{\text{no } l_{ik} \text{ or } \beta_k}
\]
Therefore, $z_i^{-k}$ is \emph{unchanged} when the L update of factor $k$
modifies $\bar{l}_{ik}$ (since only $k$-indexed quantities change).  This
means the \emph{same} $z_i^{-k}$ computed before the L update of factor $k$
can be reused for the $\beta_k$ update of factor $k$, saving computation.
In the Gauss-Seidel CAVI loop: \texttt{z\_no\_k} is computed once per factor
$k$ and reused in both step (a) [L update] and step (c) [$\beta$ update].
\end{verifybox}

The decomposition $z_i \approx z_i^{-k} + \bar{l}_{ik}\beta_k$ follows
immediately from the definition of the working response.  In code:
\begin{verbatim}
eta_no_k <- as.vector(EL %*% EBeta) - EL[, k] * EBeta[k]
z_no_k   <- z - eta_no_k   # = z_i - sum_{k'!=k} l_bar * beta_bar
\end{verbatim}

% =============================================================================
\section{Expanding the Survival Likelihood Over $q(L)$}\label{sec:expand}
% =============================================================================

Substituting the working-response form~\eqref{eq:taylor-result} into the
objective~\eqref{eq:objective-beta}:
\begin{equation}
  \mathcal{F}(q_\beta, g_\beta) \approx
  \mathbb{E}_{q_l, q_\beta}\!\left[\sum_i -\frac{1}{2}W_{ii}\!\left(z_i - \sum_k l_{ik}\beta_k\right)^{\!2}\right]
  + \sum_k \mathbb{E}_{q_{\beta_k}}\!\left[\log\frac{g_{\beta_k}(\beta_k)}{q_{\beta_k}(\beta_k)}\right]
  \label{eq:obj-expanded}
\end{equation}

\subsection{Expanding the Squared Term}

\begin{derivbox}
\textbf{Step 4.1: Expand the square.}
\begin{align}
  \left(z_i - \sum_k l_{ik}\beta_k\right)^{\!2}
  &= z_i^2 - 2z_i\sum_k l_{ik}\beta_k + \left(\sum_k l_{ik}\beta_k\right)^{\!2}
  \label{eq:expand-sq}
\end{align}

\textbf{Step 4.2: Expand the cross-product term.}
\begin{align}
  \left(\sum_k l_{ik}\beta_k\right)^{\!2}
  &= \sum_k l_{ik}^2\beta_k^2 + 2\sum_{k < k'} l_{ik}\beta_k l_{ik'}\beta_{k'}
  \label{eq:expand-cross}
\end{align}

\begin{correctionbox}
\textbf{Correction R5 (from REVISED.tex).} The document
\texttt{MF\_Derivations\_UpdateAlgo\_2\_12\_26.pdf} (Eq.~68) and also the
March~6, 2026 draft (Eq.~68 on p.~7) contain a typographical error where
the exponent on $\beta_k$ is missing: $\sum_k l^2_{ik}\beta_k$ should be
$\sum_k l^2_{ik}\beta^2_k$.  The \emph{final} formulas for $A_k$ and $B_k$
in those documents are correct; only this intermediate step has the typo.
\end{correctionbox}

\textbf{Step 4.3: Apply $\mathbb{E}_{q_l}[\,\cdot\,]$ using mean-field independence.}

Under mean-field variational inference, $q(L) = \prod_k q_{l_k}(\bm{l}_{\cdot k})$
and $l_{ik} \perp l_{ik'}$ for $k \neq k'$ under $q$.  Therefore:
\begin{align*}
  \mathbb{E}_{q_l}[l_{ik}] &= \bar{l}_{ik} \\
  \mathbb{E}_{q_l}[l_{ik}^2] &= \overline{l^2_{ik}} = \mathrm{Var}_q(l_{ik}) + \bar{l}_{ik}^2 \\
  \mathbb{E}_{q_l}[l_{ik} l_{ik'}] &= \bar{l}_{ik}\,\bar{l}_{ik'} \quad (k \neq k')
\end{align*}

The cross terms $\mathbb{E}[l_{ik}\beta_k l_{ik'}\beta_{k'}] = \bar{l}_{ik}\bar{l}_{ik'}\beta_k\beta_{k'}$
for $k \neq k'$.  These cross-factor terms survive the L expectation but are
subsequently absorbed into the $z_i^{-k}$ decomposition.

\textbf{Step 4.4: Variance-mean decomposition.}

After taking $\mathbb{E}_{q_l}$ and applying the decomposition
$\overline{l^2_{ik}} = \mathrm{Var}_q(l_{ik}) + \bar{l}_{ik}^2$, the
squared term splits as:
\begin{equation}
  \mathbb{E}_{q_l}\!\left[\left(\sum_k l_{ik}\beta_k\right)^{\!2}\right]
  = \underbrace{\left(\sum_k \bar{l}_{ik}\beta_k\right)^{\!2}}_{\text{plug-in square}}
  + \underbrace{\sum_k \mathrm{Var}_q(l_{ik})\,\beta_k^2}_{\text{variance correction}}
  \label{eq:variance-split}
\end{equation}

\textbf{Step 4.5: Reconstitute the full objective.}

The expectation of the squared-residual term becomes:
\begin{align}
  \mathbb{E}_{q_l}\!\left[\left(z_i - \sum_k l_{ik}\beta_k\right)^{\!2}\right]
  &= \left(z_i - \sum_k \bar{l}_{ik}\beta_k\right)^{\!2}
     + \sum_k \mathrm{Var}_q(l_{ik})\,\beta_k^2
  \label{eq:exp-sq-result}
\end{align}

Substituting into the objective~\eqref{eq:obj-expanded}:
\begin{align}
  \mathcal{F}(q_\beta, g_\beta) &\approx
  \mathbb{E}_{q_\beta}\!\left[\sum_i -\frac{1}{2}W_{ii}\left(z_i - \sum_k \bar{l}_{ik}\beta_k\right)^{\!2}\right]
  \nonumber \\
  &\quad + \mathbb{E}_{q_\beta}\!\left[\sum_i -\frac{1}{2}W_{ii}\sum_k \mathrm{Var}_q(l_{ik})\,\beta_k^2\right]
  \nonumber \\
  &\quad + \sum_k \mathbb{E}_{q_{\beta_k}}\!\left[\log\frac{g_{\beta_k}(\beta_k)}{q_{\beta_k}(\beta_k)}\right]
  \label{eq:objective-full}
\end{align}
\end{derivbox}

% =============================================================================
\section{Coordinate-Ascent for a Single $\beta_k$}\label{sec:coordinate}
% =============================================================================

Apply the mean-field independence $q(\bm{\beta}) = \prod_k q_{\beta_k}(\beta_k)$
and fix all $\beta_{k'}$ for $k' \neq k$.  Substitute
$z_i = z_i^{-k} + \bar{l}_{ik}\beta_k$ into the first term of
\eqref{eq:objective-full}:

\begin{equation}
  \sum_i -\frac{1}{2}W_{ii}\left(z_i - \sum_{k} \bar{l}_{ik}\beta_k\right)^{\!2}
  = \sum_i -\frac{1}{2}W_{ii}\left(z_i^{-k} - \bar{l}_{ik}\beta_k\right)^{\!2}
  + C
  \label{eq:ca-step1}
\end{equation}

where $C$ collects all terms constant in $\beta_k$.  Expanding the square:
\begin{equation}
  -\frac{1}{2}W_{ii}\!\left(z_i^{-k} - \bar{l}_{ik}\beta_k\right)^{\!2}
  = -\frac{1}{2}W_{ii}\!\left[\!\left(z_i^{-k}\right)^{\!2}
    - 2z_i^{-k}\bar{l}_{ik}\beta_k + \bar{l}_{ik}^2\beta_k^2\right]
  \label{eq:ca-expand}
\end{equation}

Summing over $i$, summing over the variance-correction term, and collecting
the $\beta_k^2$ and $\beta_k$ coefficients:

\begin{derivbox}
\textbf{Identifying the quadratic form.}

The objective for a single $\beta_k$ (up to constants in $\beta_k$) is:
\begin{align}
  \mathcal{F}(q_{\beta_k}, g_{\beta_k}) &=
  \mathbb{E}_{q_{\beta_k}}\!\Biggl[
    \underbrace{\sum_i W_{ii} z_i^{-k} \bar{l}_{ik}}_{B_k}\,\beta_k
    - \frac{1}{2}\underbrace{\sum_i W_{ii}\overline{l^2_{ik}}}_{A_k}\,\beta_k^2
  \Biggr]
  + \mathbb{E}_{q_{\beta_k}}\!\left[\log\frac{g_{\beta_k}(\beta_k)}{q_{\beta_k}(\beta_k)}\right]
  + C
  \label{eq:quadratic-form}
\end{align}

where we have combined the plug-in squared-mean term $\bar{l}_{ik}^2$ with
the variance correction $\mathrm{Var}_q(l_{ik})$ into the full second moment
$\overline{l^2_{ik}} = \bar{l}_{ik}^2 + \mathrm{Var}_q(l_{ik})$:
\begin{align}
  \bar{l}_{ik}^2 + \mathrm{Var}_q(l_{ik}) &= \overline{l^2_{ik}}
  \label{eq:second-moment-combine}
\end{align}

This yields the precision and signal coefficients:
\begin{align}
  A_k &= \sum_{i=1}^n W_{ii}\,\overline{l^2_{ik}}
  \label{eq:Ak} \\
  B_k &= \sum_{i=1}^n W_{ii}\,z_i^{-k}\,\bar{l}_{ik}
  \label{eq:Bk}
\end{align}
\end{derivbox}

\begin{verifybox}
\textbf{Sign and dimension checks.}
\begin{itemize}
  \item $A_k > 0$: since $W_{ii} > 0$ (floored at $10^{-6}$) and
        $\overline{l^2_{ik}} \geq 0$, we have $A_k \geq 0$.  A precision
        floor $A_k \geq 10^{-10}$ is applied in code to prevent division
        by zero.
  \item $A_k$ is a \textbf{scalar} (sum over $i = 1,\ldots,n$).
  \item $B_k$ is a \textbf{scalar} (sum over $i = 1,\ldots,n$).
  \item The quadratic $-\frac{1}{2}A_k\beta_k^2 + B_k\beta_k$ is concave in
        $\beta_k$ (since $A_k > 0$), as required for a valid EBNM likelihood.
\end{itemize}
\end{verifybox}

% =============================================================================
\section{The EBNM Problem}\label{sec:ebnm}
% =============================================================================

The objective~\eqref{eq:quadratic-form} is exactly a one-dimensional EBNM
problem: maximise over $(q_{\beta_k}, g_{\beta_k})$ the objective
$\mathbb{E}_q[-\frac{1}{2A_k}(\beta_k - x_k)^2] + \mathbb{E}_q[\log(g/q)]$
where:

\begin{ebnmbox}
\textbf{EBNM Inputs:}
\begin{align}
  x_k &= \frac{B_k}{A_k}
  = \frac{\displaystyle\sum_i W_{ii}\,z_i^{-k}\,\bar{l}_{ik}}
         {\displaystyle\sum_i W_{ii}\,\overline{l^2_{ik}}}
  \label{eq:xk} \\[6pt]
  s_k &= \frac{1}{\sqrt{A_k}}
  = \frac{1}{\sqrt{\displaystyle\sum_i W_{ii}\,\overline{l^2_{ik}}}}
  \label{eq:sk}
\end{align}

\textbf{EBNM Call:}
\[
  (\hat{g}_{\beta_k},\,\hat{q}_{\beta_k}) = \mathrm{EBNM}(x_k,\, s_k)
\]

\textbf{Posterior Extraction:}
\begin{align}
  \bar{\beta}_k &= \hat{q}_{\beta_k}\text{.posterior\$mean}
  \label{eq:beta-mean} \\
  \overline{\beta^2_k} &= \hat{q}_{\beta_k}\text{.posterior\$sd}^2
                          + \hat{q}_{\beta_k}\text{.posterior\$mean}^2
  \label{eq:beta-sec}
\end{align}
\end{ebnmbox}

\subsection{Prior Variance Note}

The derivation with a fixed Gaussian prior $g(\beta_k) \sim \mathcal{N}(0, \sigma^2)$
would yield an additional term $1/\sigma^2$ in $A_k$:
\begin{equation}
  A_k^{\text{fixed prior}} = \sum_i W_{ii}\overline{l^2_{ik}} + \frac{1}{\sigma^2}
  \label{eq:Ak-with-prior}
\end{equation}

In the implementation, $\sigma^2$ is \emph{not} fixed in advance.  Instead,
\texttt{prior\_family = "point\_normal"} causes the EBNM solver to
\emph{empirically estimate} the prior from data (empirical Bayes).  This is
equivalent to marginalising over $\sigma^2$, and the $1/\sigma^2$ term is
absorbed into the EBNM prior estimation.  See REVISED.tex lines 911--914.

\subsection{Error-in-Variables Interpretation}

\begin{verifybox}
\textbf{Error-in-variables correction.}
The use of $\overline{l^2_{ik}}$ (full second moment) rather than
$\bar{l}_{ik}^2$ (squared mean) in $A_k$ encodes an
\emph{error-in-variables correction}:

\begin{enumerate}
  \item Higher posterior variance $\mathrm{Var}_q(l_{ik})$ increases
        $\overline{l^2_{ik}} = \bar{l}_{ik}^2 + \mathrm{Var}_q(l_{ik})$.
  \item This increases $A_k$, which decreases both $x_k = B_k/A_k$ and
        $s_k = 1/\sqrt{A_k}$.  Crucially, the ratio
        $|x_k|/s_k = |B_k|/\sqrt{A_k}$ decreases.
  \item A smaller $|x_k|/s_k$ ratio means the signal is weaker relative to
        the noise, so the EBNM prior applies more shrinkage.
  \item Net effect: more uncertainty in $L$ $\Rightarrow$ more shrinkage of
        $\beta_k$ toward zero.  This prevents overfitting $\beta$ to
        uncertain loadings.
\end{enumerate}

Mathematically: if $\mathrm{Var}_q(l_{ik}) = 0$ for all $i$, then
$A_k = \sum_i W_{ii}\bar{l}_{ik}^2$ and the EBNM with a flat prior
recovers the weighted least squares estimator:
\[
  \hat{\beta}_k^{\mathrm{WLS}} = \frac{\sum_i W_{ii} z_i^{-k} \bar{l}_{ik}}
                                       {\sum_i W_{ii} \bar{l}_{ik}^2}
\]
This is verified in test T2.1 of \texttt{tests/test\_update\_beta.R}.
\end{verifybox}

% =============================================================================
\section{Verification and Self-Consistency Checks}\label{sec:verify}
% =============================================================================

\begin{verifybox}
\textbf{Check 1: WLS limit.}
When posterior variance is zero ($\overline{l^2_{ik}} = \bar{l}_{ik}^2$)
and the EBNM prior is flat (non-informative), the posterior mean equals the
WLS estimator $x_k = B_k/A_k$.  \textbf{Confirmed in test T2.1.}
\end{verifybox}

\begin{verifybox}
\textbf{Check 2: Null factor shrinkage.}
When $z_i^{-k}$ is pure noise (no signal from factor $k$), $B_k \approx 0$
and $x_k \approx 0$.  The point-normal prior shrinks $\bar{\beta}_k$ to
approximately zero.  \textbf{Confirmed in tests T5.1--T5.2.}
\end{verifybox}

\begin{verifybox}
\textbf{Check 3: Error-in-variables direction.}
Larger $\overline{l^2_{ik}}$ (more posterior variance in $L$) yields larger
$A_k$ and a smaller ratio $|B_k/A_k|/s_k = |B_k|/\sqrt{A_k}$, leading to
more shrinkage.  Note: $s_k = 1/\sqrt{A_k}$ is \emph{smaller} (not larger)
when $A_k$ is larger.  The shrinkage effect comes from the reduction of the
pseudo-observation $x_k$, not from an increase in $s_k$.
\textbf{Confirmed in test T6.1.}
\end{verifybox}

\begin{verifybox}
\textbf{Check 4: Consistency with V2.R.}
The modular function \texttt{update\_beta\_k()} produces identical results
(to machine precision) as the inline code in
\texttt{Supervised\_Bayesian\_MF\_V2.R} lines 356--361 for all tested inputs.
\textbf{Confirmed in test T9.1.}
\end{verifybox}

\begin{verifybox}
\textbf{Check 5: z\_no\_k reuse invariance.}
$z_i^{-k}$ does not depend on $l_{ik}$ or $\beta_k$.  Therefore the value
computed before the L update of factor $k$ remains valid for use in the
$\beta_k$ update.  \textbf{Confirmed analytically and in test T8.1.}
\end{verifybox}

% =============================================================================
\section{Code Mapping}\label{sec:code}
% =============================================================================

\begin{codebox}
\textbf{Math $\to$ Code correspondence for \texttt{update\_beta.R}:}

\begin{center}
\begin{tabular}{lll}
\toprule
\textbf{Math} & \textbf{R Variable} & \textbf{Computation} \\
\midrule
$W_{ii}$ & \texttt{w} & \texttt{calc\_cox\_taylor(eta,time,status)\$w} \\
$z_i$ & \texttt{z} & \texttt{eta + u/w} \\
$z_i^{-k}$ & \texttt{z\_no\_k} & \texttt{compute\_z\_no\_k(z, EL, EBeta, k)} \\
$\bar{l}_{ik}$ & \texttt{EL[i,k]} & \texttt{res\_L\$posterior\$mean} \\
$\overline{l^2_{ik}}$ & \texttt{EL2[i,k]} & \texttt{sd\^{}2 + mean\^{}2} \\
$A_k$ & \texttt{res\$A} & \texttt{max(sum(w * EL2\_k), 1e-10)} \\
$B_k$ & \texttt{res\$B} & \texttt{sum(w * z\_no\_k * EL\_k)} \\
$x_k$ & \texttt{res\$x} & \texttt{B\_k / A\_k} \\
$s_k$ & \texttt{res\$s} & \texttt{1 / sqrt(A\_k)} \\
$\bar{\beta}_k$ & \texttt{res\$mean} & \texttt{ebnm\_result\$posterior\$mean} \\
$\overline{\beta^2_k}$ & \texttt{res\$second} & \texttt{sd\^{}2 + mean\^{}2} \\
\bottomrule
\end{tabular}
\end{center}

\textbf{Pseudocode} matching V2.R lines 356--361:
\begin{verbatim}
A_k   <- max(sum(w * EL2[, k]), 1e-10)          # precision
B_k   <- sum(w * z_no_k * EL[, k])              # signal
x_k   <- B_k / A_k                              # pseudo-obs
s_k   <- 1 / sqrt(A_k)                          # pseudo-noise
res   <- ebnm(x = x_k, s = s_k,
              prior_family = "point_normal")
EBeta[k]  <- res$posterior$mean                 # posterior mean
EBeta2[k] <- res$posterior$sd^2 + mean^2        # 2nd moment
\end{verbatim}
\end{codebox}

% =============================================================================
\section{Algorithm Summary}\label{sec:summary}
% =============================================================================

The complete $q_{\beta_k}$ update algorithm (Gauss-Seidel over $k = 1,\ldots,K$):

\begin{enumerate}
  \item \textbf{Prerequisite:} Compute Cox Taylor quantities $(z_i, W_{ii})$ from
        \texttt{calc\_cox\_taylor}$(\hat{\bm{\eta}},\, t,\, \delta)$.
  \item \textbf{For} $k = 1, \ldots, K$:
  \begin{enumerate}
    \item Compute partial working response:
          $z_i^{-k} = z_i - \sum_{k'\neq k}\bar{l}_{ik'}\bar{\beta}_{k'}$
    \item Compute precision: $A_k = \max\!\left(\sum_i W_{ii}\overline{l^2_{ik}},\,10^{-10}\right)$
    \item Compute signal: $B_k = \sum_i W_{ii} z_i^{-k} \bar{l}_{ik}$
    \item Form EBNM inputs: $x_k = B_k/A_k$, $s_k = 1/\sqrt{A_k}$
    \item Solve: $(\hat{g}_{\beta_k}, \hat{q}_{\beta_k}) = \mathrm{EBNM}(x_k, s_k)$
    \item Extract: $\bar{\beta}_k \leftarrow \hat{q}$.posterior\$mean;
          $\overline{\beta^2_k} \leftarrow \hat{q}$.posterior\$sd$^2$ + mean$^2$
    \item \textbf{Gauss-Seidel:} use updated $\bar{\beta}_k$ when computing
          $z_i^{-k'}$ for subsequent $k' > k$
  \end{enumerate}
\end{enumerate}

% =============================================================================
\appendix
\section{Notation Index}\label{app:notation}
% =============================================================================

\begin{center}
\begin{tabular}{lll}
\toprule
\textbf{Symbol} & \textbf{Dimension} & \textbf{Meaning} \\
\midrule
$n$ & scalar & number of patients \\
$p$ & scalar & number of genomic features \\
$K$ & scalar & number of latent factors \\
$Y$ & $n\times p$ & observed genomics data \\
$L$ & $n\times K$ & patient loading matrix \\
$F$ & $p\times K$ & biological factor matrix \\
$\bm{\beta}$ & $K\times 1$ & survival coefficient vector \\
$\tau_j$ & $p$-vector & feature-specific noise precision \\
$(t_i, \delta_i)$ & $n$-vectors & survival time and event indicator \\
$\eta_i$ & $n$-vector & linear predictor $\bm{l}_i^\top\bm{\beta}$ \\
$u_i$ & $n$-vector & Cox score \\
$W_{ii}$ & $n$-vector & Cox neg-diagonal Hessian (positive) \\
$z_i$ & $n$-vector & working response $\hat{\eta}_i + u_i/W_{ii}$ \\
$z_i^{-k}$ & $n$-vector & partial working response (excludes factor $k$) \\
$A_k$ & scalar & EBNM precision for factor $k$ \\
$B_k$ & scalar & EBNM signal for factor $k$ \\
$x_k$ & scalar & EBNM pseudo-observation $B_k/A_k$ \\
$s_k$ & scalar & EBNM pseudo-noise $1/\sqrt{A_k}$ \\
\bottomrule
\end{tabular}
\end{center}

% =============================================================================
\section{Cross-References to REVISED.tex}\label{app:xref}
% =============================================================================

\begin{center}
\begin{tabular}{ll}
\toprule
\textbf{This document} & \textbf{REVISED.tex} \\
\midrule
Eq.~\eqref{eq:genomics}--\eqref{eq:survival} & Sec. 1A (model specification) \\
Eq.~\eqref{eq:elbo-full} & Eq.~(1) (full ELBO) \\
Eq.~\eqref{eq:score}--\eqref{eq:working-response} & Sec. 3 (Taylor expansion) \\
Eq.~\eqref{eq:Ak}--\eqref{eq:Bk} & Eqs.~(79)--(80) ($A_k$, $B_k$) \\
Eq.~\eqref{eq:xk}--\eqref{eq:sk} & Eqs.~(81)--(82) (EBNM inputs) \\
Correction R5 & Sec. 1 (revision log) \\
\bottomrule
\end{tabular}
\end{center}

\end{document}
